\subsection{商环的维数}

命题 1. --- 设 A 是一个诺特整环，  是 A 的非零元素，  是包含  的 A 的素理想集合中的一个极小元。则  的高度为 1。

设 \(\mathfrak{q} \subset \mathfrak{p}\)；\(\mathfrak{p}\) 是一个素理想，满足 \(x \notin \mathfrak{q}\)，且不同于 \(\mathfrak{p}\)。由 \(\mathrm{A}_{\mathfrak{p}}\) 的极小性，\(\mathrm{A}_{\mathfrak{q}}\) 是一个局部环。由于 A 是整环，对每个整数 \(n \geqslant 0\)，记 \(\mathfrak{q}_n\) 为 \(\mathfrak{q}^n\mathrm{A}_{\mathfrak{q}} \cap \mathrm{A}_{\mathfrak{p}}\)；这是 \(\mathrm{A}_{\mathfrak{p}}\) 的一个理想。\(\mathfrak{p}\) 的极小性意味着局部环 \(\mathrm{A}_{\mathfrak{p}} / x\mathrm{A}_{\mathfrak{p}}\) 的维数为 0；因而它具有有限长度（§ 1, no. 3, 例 1），存在整数 \(n_0 \geqslant 0\)，使得

\[
\mathfrak {q} _ {n} + x \mathrm{A} _ {\mathfrak {p}} = \mathfrak {q} _ {n + 1} + x \mathrm{A} _ {\mathfrak {p}} \quad \text { for   tout } \quad n \geqslant n _ {0}.\tag{1}
\]

固定整数 \(n \geqslant n_{0}\)。给定 \(y \in \mathfrak{q}_{n}\)，存在 \(a \in \mathrm{A}_{\mathfrak{p}}\)，使 \(y - ax \in \mathfrak{q}_{n+1}\)；于是 \(ax \in \mathfrak{q}_{n}\)，从而 \(a \in \mathfrak{q}_{n}\)，且 \(x \notin \mathfrak{q}\)，最后有 \(y \in \mathfrak{q}_{n+1} + x\mathfrak{q}_{n}\)。因此有

\[
\mathfrak {q} _ {n} = \mathfrak {q} _ {n + 1} + x \mathfrak {q} _ {n}.\tag{2}
\]

由于  属于诺特局部环  的极大理想，中山引理表明 \(A_{p}\)\(\mathfrak{q}_n = \mathfrak{q}_{n+1}\)。由于 \((\mathfrak{q}A_{\mathfrak{q}})^{n} = \mathfrak{q}_{n}A_{\mathfrak{q}}\)，因此得到

\[
(q A _ {q}) ^ {n} = (q A _ {q}) ^ {n + 1} \quad \text { for   tout } \quad n \geqslant n _ {0}.\tag{3}
\]

由于环 \(\mathbf{A}_{\mathfrak{q}}\) 是局部诺特环，有 \(\bigcap_{n\geqslant 0}(\mathfrak{q}\mathbf{A}_{\mathfrak{q}})^n = \{0\}\)（III, § 3, no. 2, 命题 5 的推论），从而 \((\mathfrak{q}\mathbf{A}_{\mathfrak{q}})^{n_0} = \{0\}\)，最后，\(\mathfrak{q}\mathbf{A}_{\mathfrak{q}}\) 这个 \(\mathbf{A}_{\mathfrak{q}}\) 的素理想退化为 0。因此 \(\mathfrak{q} = \{0\}\)，这证明了 \(\mathfrak{p}\) 的高度为 1。

命题 2. --- 设 A 是一个诺特环，  是一个正整数，  是包含于 A 的根基中且由  个元素生成的理想。则有

\[
\dim (\mathrm{A} / \mathfrak {a}) \leqslant \dim (\mathrm{A}) \leqslant \dim (\mathrm{A} / \mathfrak {a}) + m.\tag{4}
\]

不等式  由 § 1, no. 3 的命题 6 得出。对  作直接归纳可知，只需对 A 的根基中的每个元素 x 证明不等式

\[
\dim (\mathrm{A}) \leqslant \dim (\mathrm{A} / x \mathrm{A}) + 1\tag{5}
\]

对 A 的根基中的每个元素 \(x\)，即要证明 \(\dim(A/xA) \geqslant n - 1\) 对于 A 的每条素理想链 \(\mathfrak{p}_0 \subset \ldots \subset \mathfrak{p}_n\)，其长度为 \(n \geqslant 1\)，且 \(x \in \mathfrak{p}_n\)。只需构造一条 A 的素理想链 \(\mathfrak{q}_1 \subset \ldots \subset \mathfrak{q}_n\)，使 \(x \in \mathfrak{q}_1\)；这由下述引理得出：

引理 1. --- 设 A 是一个诺特环，\(\mathfrak{p}_{0} \subset \ldots \subset \mathfrak{p}_{n}\) 是 A 的一条素理想链，其长度为 \(n \geqslant 1\)，且 \(x\) 属于 \(\mathfrak{p}_{n}\)。存在一条 \(\mathfrak{p}_{0}' \subset \ldots \subset \mathfrak{p}_{n}'\) 的链，使 \(\mathfrak{p}_{0}' = \mathfrak{p}_{0}\)、\(\mathfrak{p}_{n}' = \mathfrak{p}_{n}\)，且 \(x \in \mathfrak{p}_{1}'\)。

我们对 \(n\) 作归纳，\(n=1\) 时显然成立。因此设 \(n\geqslant2\)，且 \(x\) 不属于 \(\mathfrak{p}_{n-1}\)。令 \(\mathfrak{p}_{n-1}^{\prime}\) 为包含于 \(\mathfrak{p}_{n}^{\prime}=\mathfrak{p}_{n}\) 且包含 \(\mathfrak{p}_{n-2}+\mathbf{A}x\) 的 A 的素理想集合中的一个极小元（II, § 2, no. 6, 引理 2）。根据命题 1，\(\mathfrak{p}_{n-1}^{\prime}/\mathfrak{p}_{n-2}\) 这个环 \(\mathbf{A}/\mathfrak{p}_{n-2}\) 的理想具有高度 1；又由于 \(\mathfrak{p}_{n-2}\subset\mathfrak{p}_{n-1}\subset\mathfrak{p}_{n}\) 是长度为 2 的链，对于 \(\mathfrak{p}_{n-2}\subset\mathfrak{p}_{n-1}^{\prime}\subset\mathfrak{p}_{n}^{\prime}\) 也同样如此。有 \(x\in\mathfrak{p}_{n-1}^{\prime}\)。将归纳假设应用于链 \(\mathfrak{p}_{0}\subset\mathfrak{p}_{1}\subset\ldots\subset\mathfrak{p}_{n-2}\subset\mathfrak{p}_{n-1}^{\prime}\)，可知存在一条 \(\mathfrak{p}_{0}^{\prime}\subset\mathfrak{p}_{1}^{\prime}\subset\ldots\subset\mathfrak{p}_{n-2}^{\prime}\subset\mathfrak{p}_{n-1}^{\prime}\) 的链，使 \(x\in\mathfrak{p}_{1}^{\prime}\) 且 \(\mathfrak{p}_{0}^{\prime}=\mathfrak{p}_{0}\)。链

\[
\mathfrak {p} _ {0} ^ {\prime} \subset \mathfrak {p} _ {1} ^ {\prime} \subset \dots \subset \mathfrak {p} _ {n - 1} ^ {\prime} \subset \mathfrak {p} _ {n} ^ {\prime}
\]

满足所要求的条件。

推论 1. --- a) 每个诺特局部环都有有限维数。更一般地，每个非零诺特半局部环（II, § 3, no. 5, 定义 4）都有有限维数。

\begin{enumerate}
\def\labelenumi{\alph{enumi})}
\setcounter{enumi}{1}
\item
  设 A 是一个诺特环。A 的每个异于 A 的理想都有有限高度。
\item
  诺特环 A 的每个素理想递减序列都是驻定的。
\item
  设 A 是一个非零诺特半局部环，  是它的根基；商环  是非零的阿廷环，因而维数为 0（§ 1, no. 3, 例 1）。存在一个整数  ，使得 A 的理想  由  个元素生成；因此根据命题 2 有  。
\item
  设  是 A 的一个理想，且令 \(\mathfrak{a} \neq A\) 是包含 \(\mathfrak{a}\) 的 A 的一个极大理想。根据 § 1, no. 3 的命题 7，有  ，并且 \(0 \leqslant \mathrm{ht}(\mathfrak{a}) \leqslant \dim(A_{\mathfrak{m}})\)\(A_{\mathfrak{m}}\) 是诺特局部环。因此由 a)，\(\mathrm{ht}(\mathfrak{a})\) 有限。
\item
  A 的每个有限严格递减素理想序列 \((\mathfrak{p}_{i})_{0\leqslant i\leqslant n}\) 都定义了一条链  ，从而 \(\mathfrak p_{n}\subset\ldots\subset\mathfrak p_{0}\)\(n\leqslant\dim(A_{\mathfrak p_{0}})<+\infty\)。因此不可能存在 A 的无限严格递减素理想序列，从而得到 c)。
\end{enumerate}

推论 2. --- 设 A 是一个诺特局部环。

\begin{enumerate}
\def\labelenumi{\alph{enumi})}
\item
设 \(x \in \mathfrak{m}_{\mathbf{A}}\)。则 \(\dim(\mathbf{A}/x\mathbf{A})\) 等于 \(\dim(\mathbf{A})\) 或 \(\dim(\mathbf{A}) - 1\)。要使 \(\dim(\mathbf{A}/x\mathbf{A}) = \dim(\mathbf{A}) - 1\)，必要且充分条件是：\(x\) 不属于任何极小素理想 \(\mathfrak{p}\)，这些理想属于 \(\mathbf{A}\) 且满足 \(\dim(\mathbf{A}/\mathfrak{p}) = \dim(\mathbf{A})\)，并且只要 \(x\) 不是 \(\mathbf{A}\) 中的零因子就足够。
\item
  设 \(\mathfrak{a}\) 是 A 的一个异于 A 的理想，且满足 \(\dim(\mathbf{A}/\mathfrak{a}) < \dim(\mathbf{A})\)。存在 \(x \in \mathfrak{a}\)，使得 \(\dim(\mathbf{A}/x\mathbf{A}) = \dim(\mathbf{A}) - 1\)。
\item
  如果  ，则存在  ，使得 \(	ext{If} \dim(A) \geqslant 1\)\(there exists x \in \mathfrak{m}_{A} \text{ \text{使得} } \dim(A/xA) = \dim(A) - 1\)。
\end{enumerate}

根据命题 2，  等于  或  。要使 \(\dim(A/xA) = \dim(A)\) 成立，必要且充分条件是存在 A 的一条素理想链 \(\mathfrak{p}_0 \subset \ldots \subset \mathfrak{p}_n\)，使得 \(x \in \mathfrak{p}_0\) 且 \(n = \dim(A)\)，也就是说，存在 A 的一个素理想 \(\mathfrak{p}_0\)，包含 \(x\)，且满足 \(\dim(A/\mathfrak{p}_0) = \dim(A)\)。但是，这样的素理想 \(\mathfrak{p}_0\) 必然是极小的，因而 \(\mathfrak{p}_0\) 的每个元素都是 A 中的零因子（IV, § 1, no. 1, 命题 2 的推论 3 以及 no. 4, 定理 2）。这证明了 \(a)\)。

令 \(\Phi\) 为 A 的极小素理想集合，令 \(\Phi'\) 为满足 \(\mathfrak{p} \in \Phi\) 且 \(\dim(A/\mathfrak{p}) = \dim(A)\) 的集合。我们知道 \(\Phi\) 是有限的（II, § 4, no. 3, 命题 14 的推论 3），所以 \(\Phi'\) 也是有限的。设 \(\mathfrak{a}\) 是 A 的一个理想，满足 \(\dim(A/\mathfrak{a}) < \dim(A)\)。对于每个 \(\mathfrak{p} \in \Phi'\)，有 \(\dim(A/\mathfrak{a}) < \dim(A/\mathfrak{p})\)，从而 \(\mathfrak{a} \not\subsetneq \mathfrak{p}\)。根据 II, § 1, no. 1 的命题 2，存在一个元素 \(x\)，属于 \(\mathfrak{a}\) 且不属于任何 \(\mathfrak{p} \in \Phi'\)，于是有 \(\dim(A/xA) = \dim(A) - 1\) 根据 \(a)\)，这证明了 \(b)\)。

断言 c) 是在 \(\mathfrak{a} = \mathfrak{m}_{\mathrm{A}}\) 时的 \(b\) 的特例。
% semantic-fragments: legacy-input-seam
\relax{}
